\section{Actual arm quotients and signed compensation of a private deep tail}
\label{sac-section}

This section keeps the negative radical contributions in the signed
large-prime cost. It proves a sufficient subclass in which two actual
quotient arms compensate for unrestricted depths in the third. General
membership in that subclass remains unproved.

Let $\mathcal O=\mathbb Z[\zeta]$, where $\zeta^2-\zeta+1=0$, and let
$w=a+b\zeta$ be primitive and unramified, with
\[
 Q=N(w)\ge7,\qquad \gcd(a,b)=1,\qquad 3\nmid Q.
\]
The three input arms $a,b,a+b$ are nonzero. For an integer $n\ge5$
with $\gcd(n,6)=1$, define
\begin{equation}\label{sac-normalized}
 z_n=A+B\zeta=
 \begin{cases}w^n,&n\equiv1\pmod6,\\
 \bar w^{\,n},&n\equiv5\pmod6,
 \end{cases}
 \qquad
 D_1=\frac A a,\quad D_2=\frac B b,\quad
 D_3=\frac{A+B}{a+b}.
\end{equation}
Write $P(x+y\zeta)=xy(x+y)$ and $T_j=|P(w^j)|$.

\begin{lemma}[Three actual integer quotients]\label{sac-quotients}
The $D_i$ are evaluations of integer homogeneous polynomials of degree
$n-1$. They are nonzero and pairwise coprime, and satisfy
\begin{equation}\label{sac-relations}
 aD_1+bD_2=(a+b)D_3,\qquad
 |D_1D_2D_3|=\frac{T_n}{T_1}.
\end{equation}
For prime $n$, their product is the actual top cyclotomic factor $C_n$
up to sign. For general $n$ in this domain it is, up to sign,
$\prod_{d\mid n,\ d>1}C_d$.
\end{lemma}
\begin{proof}
The normalized power map in \eqref{sac-normalized} fixes each boundary
line individually. At $a=0$ its value is $b^n\zeta$, at $b=0$ it is
$a^n$, and at $a=-b$ it is $b^n\zeta^2$. This uses $n\equiv1\pmod6$
in the first branch and $-n\equiv1\pmod6$ in the second.
Consequently its first coordinate is divisible by $a$ in $\mathbb Z[a,b]$,
its second by $b$, and the sum by $a+b$. For the last assertion,
substitute $a=-b$ and divide by the monic linear polynomial $a+b$.
Homogeneity gives the degrees.

The established oriented factorization shows that the unramified
primitive power $w^n$ is primitive and has no zero boundary arm.
Conjugation preserves these properties. Thus $A,B,A+B$ are pairwise
coprime nonzero integers, and their respective integer divisors $D_i$
are pairwise coprime. Substitution and multiplication give
\eqref{sac-relations}. Finally $P(\bar z)=-P(z)$ and the actual
cyclotomic identity of Theorem~\ref{cp-integers} give the last statements.
\end{proof}

The unramified hypothesis matters: for $w=1+\zeta$ and $n=5$ the
normalized output is $-9-9\zeta$, and all three quotients are $-9$.
It is excluded by $3\nmid Q$.

Set
\[
 c=\max\{|A|,|B|,|A+B|\},\quad t=\log c,\quad
 c_1=\max\{|a|,|b|,|a+b|\},\quad \Delta=3t-\log T_n.
\]
For a real cutoff $Y\ge5$, all sums below are over primes and valuations
of signed integers mean valuations of their absolute values. Define
\begin{align*}
 L(Y)&=\sum_{q\le Y}v_q(T_n)\log q,&
 S_i(Y)&=\sum_{q>Y}v_q(D_i)\log q,\\
 R_i(Y)&=\sum_{\substack{q>Y\\q\mid D_i}}\log q,&
 E_i(Y)&=\sum_{q>Y}(v_q(D_i)-2)_+\log q.
\end{align*}
As before,
\[
 \log W_Y(H)=\sum_{\substack{q>Y\\q\mid H}}(v_q(H)-3)\log q
 \qquad(H\ne0).
\]

\begin{theorem}[Finite signed-credit compensation]\label{sac-finite}
For every choice of distinct $i,j,k\in\{1,2,3\}$,
\begin{equation}\label{sac-finite-bound}
 \log W_Y(T_n)\le
 \Delta+\log c_1+\log T_1+\frac{L(Y)}2
       +\frac32(E_i(Y)+E_j(Y))-3R_k(Y).
\end{equation}
The third quotient's entire multiplicity is included; no individual
depth bound is imposed on it.
\end{theorem}
\begin{proof}
Each output arm has absolute value at most $c$. Since the sum of their
logarithms is $3t-\Delta$, each has logarithm at least $t-\Delta$.
The nonzero integer input arms have absolute values between one and
$c_1$. Therefore, for $M_i=\log|D_i|$,
\begin{equation}\label{sac-arm-height}
 t-\Delta-\log c_1\le M_i\le t.
\end{equation}
Let $L_i=M_i-S_i$ be the small-prime mass of $D_i$.
Equation~\eqref{sac-relations} gives $L_i+L_j\le L(Y)$.
At every supported prime, $e\le2+(e-2)_+$, so
$2R_i\ge S_i-E_i$ and $2R_j\ge S_j-E_j$.
Pairwise coprimality makes the signed cost additive on the quotient
product. It follows that
\[
 \log W_Y(D_1D_2D_3)
 \le S_k-\frac{S_i+S_j}{2}+\frac32(E_i+E_j)-3R_k.
\]
By \eqref{sac-arm-height}, the first two terms on the right are at most
$\Delta+\log c_1+(L_i+L_j)/2$.

For arbitrary nonzero integers $U,V$, adjoining the valuations of $U$
can increase $\log W_Y(V)$ by at most $\log|U|$. At an old prime the
increase is exactly $v_q(U)\log q$; at a new prime it is
$(v_q(U)-3)\log q\le v_q(U)\log q$. Applying this to
$T_n=T_1|D_1D_2D_3|$ proves \eqref{sac-finite-bound}. This step makes
no coprimality assumption between $T_1$ and the quotient product and
does not discard a negative contribution at their overlapping primes.
\end{proof}

\begin{theorem}[Uniform explicit error and a signed-tail subclass]
\label{sac-subclass}
There are absolute effective constants $A_1,A_2$ such that, on the
whole actual domain above,
\begin{align}
 \frac{\Delta}{t}&\le A_1\frac{\log(4n)}n,&
 \frac{v_q(T_n)}t&\le A_2q^2\frac{\log^2(4n)}n
                    \quad(q\text{ prime}),\label{sac-two-place}\\
 \frac{\log c_1+\log T_1}{t}&\le\frac5n,&
 \frac{L(Y)}t&\le A_2\frac{Y^3\log Y\log^2(4n)}n.
                    \label{sac-explicit-errors}
\end{align}
Consequently, writing
\[
 \epsilon(n,Y)=A_1\frac{\log(4n)}n+\frac5n+
             \frac{A_2Y^3\log Y\log^2(4n)}{2n},
\]
one has
\begin{equation}\label{sac-uniform}
 \frac{\log W_Y(T_n)}t
 \le\epsilon(n,Y)+\frac{3}{2t}(E_i+E_j-2R_k).
\end{equation}
For $Y=n^{1/6}$, $\epsilon(n,Y)\to0$ uniformly in the moving root.
In particular, if two actual quotient arms have $v_q(D_i),v_q(D_j)\le2$
at every prime $q>Y$, then
\begin{equation}\label{sac-positive-part}
 \frac{(\log W_Y(T_n))_+}{t}\longrightarrow0
 \qquad(n\longrightarrow\infty)
\end{equation}
uniformly in that subclass. More generally it suffices that
$ (E_i+E_j-2R_k)_+/t\to0$.
\end{theorem}
\begin{proof}
To apply Theorem~\ref{lr-two-place}, multiply $z_n$ by a unit to
place it in the open positive sector. No boundary arm is zero, so such
a unit exists. Units and conjugation only permute the three absolute
arm values, preserving $c,T_n,\Delta$ and all valuations. The resulting
actual representation has ramified exponent zero, unit residual,
root $w$ or $\bar w$ of norm $Q$, and exponent $g=n$. Thus
$\lambda=1/n$, proving \eqref{sac-two-place}. This includes $q=2,3$;
primes dividing $Q$ have no boundary contribution.

The quadratic and cubic height inequalities give
$Q^n\le c^2$, $Q\ge3c_1^2/4$ and $T_1\le Q^{3/2}$.
The numerator $\log c_1+\log T_1$ is at most
$2\log Q+\log(2/\sqrt3)$. Divide by $t\ge n\log Q/2$ and use
\[
 4+\frac{2\log(2/\sqrt3)}{\log7}<5.
\]
There are at most $Y$ primes below $Y$, and each satisfies
$q^2\log q\le Y^2\log Y$. Summing \eqref{sac-two-place} proves
the remaining inequality in \eqref{sac-explicit-errors}.
Substitution in Theorem~\ref{sac-finite} gives \eqref{sac-uniform}.
For $Y=n^{1/6}$ its error tends to zero and eventually $Y\ge5$.
If the two selected arms have depth at most two, then $E_i=E_j=0$;
the term $-3R_k/t$ is nonpositive. Taking positive parts gives
\eqref{sac-positive-part}. The more general condition follows directly
from the same upper bound.
\end{proof}

Together with the previously proved archimedean and small-prime
absorption, this gives the restricted ABC estimate for this explicitly
specified homogeneous subclass, uniformly for sufficiently large $n$
for each fixed target epsilon. It does not settle bounded $n$ with
arbitrarily moving roots. Finitely many indices are not finitely many
seeds. The signed upper bound is one-sided; a persistent negative
contribution is compatible with it.

\paragraph{Exact finite evidence and formal scope.}
The exact replay checks six polynomial quotient identities, $113$
actual unramified primitive examples and $678$ rational inequalities
obtained by exponentiating twice \eqref{sac-finite-bound}. It uses
integer arithmetic, rational fractions and complete trial division,
not floating-point logarithms. For example, $w=54345+\zeta$, $n=5$ gives
\begin{align*}
 D_1&=5\cdot31^4\cdot168541\cdot11208719,\\
 D_2&=-571\cdot640049\cdot119336975731,\\
 D_3&=181\cdot659\cdot2131\cdot263129\cdot130411.
\end{align*}
Every displayed factor is prime and every product is checked exactly.
Thus two actual quotients are squarefree while the third has private
depth four at $31$. This is a finite illustration, not an asymptotic
membership theorem.

The separate Lean module \texttt{SignedArmArithmetic} proves $32$
integer declarations using the actual pair arithmetic: normalized
unit powers and congruences, the three actual divisibilities,
quotient reconstruction, the additive and boundary-product identities,
the exact norm, and quotient coprimality under an explicit output
primitivity hypothesis. It also proves the signed linear budget for
integer weights, its specialization to finite lists of natural depths
with nonnegative integer weights, and the local weighted overlap bound.
Its fresh Lean~4.32.0 build checks all $32$ declarations and $39$
dependency declarations with warnings treated as errors and only
\texttt{propext}, \texttt{Classical.choice}, \texttt{Quot.sound}.
It does not formalize preservation of primitivity, real logarithmic
weights, the full prime-valuation identities, Theorem~\ref{lr-two-place},
or general membership in the sufficient subclass.

The unresolved arithmetic problem is to constrain the actual coupled
quantity $E_i+E_j-2R_k$ using \eqref{sac-relations} and the common norm.
No claim that two arms are usually cubefree is made. The criterion
discards some shallow contributions in the two selected arms and is a
sufficient condition stronger than the general signed-tail target.
The exceptional roots of earlier averaging windows and their complete
large-prime tails remain uncontrolled.
